% Contents/pericias_e_maestrias.tex
\midsectiontitle{Perícias e Maestrias}

\section*{O Refinamento do Talento}

No universo de \textit{Overgrown}, os atributos representam o potencial bruto do aventureiro, mas são as \textbf{Perícias} que definem como esse poder é lapidado através da experiência. Elas representam o treinamento especializado, o estudo acadêmico e a memória muscular adquirida nas perigosas trilhas do Pináculo.

A maestria em uma perícia permite que o personagem realize feitos que ultrapassam a mera sorte, transformando esforço em resultados precisos.

\begin{rpgboxtips}{O Valor do Treinamento}
Uma perícia marcada como treinada (*) representa um campo onde o "instinto" não é suficiente. Tentar realizar um teste de \textbf{Manutenção} ou um \textbf{Medicina} sem o devido conhecimento pode resultar não apenas em falha, mas em consequências catastróficas para o aventureiro e seu grupo.
\end{rpgboxtips}

\noindent
\pericia{Adestramento}{Permite adestrar, acalmar ou controlar animais selvagens e montarias.}
\pericia{Arcanismo}{Define o conhecimento sobre magias, planos de existência e a manipulação da energia.}
\pericia{Atletismo}{A capacidade física de superar obstáculos, envolvendo saltar, escalar, nadar e correr longas distâncias.}
\pericia{Condução*}{A habilidade de operar veículos, desde carroças tradicionais até transportes movidos a engenharia rúnica.}
\pericia{Constituição}{A resiliência física pura; a capacidade de resistir ao cansaço, privação de sono, dor extrema e status de efeito.}
\pericia{Crime}{O domínio de atividades ilícitas, como arrombamento de fechaduras, prestidigitação e desarmar mecanismos de segurança.}
\originlink{per:chef}
\pericia{Chef*}{A arte de preparar alimentos e misturar ingredientes para criar refeições que restauram o vigor, garantem bônus temporários ou neutralizam toxinas.}
\pericia{Cultura}{O acervo de conhecimentos sobre a história do mundo, etiquetas sociais, heráldica, línguas regionais e religiões.}
\pericia{Discernimento}{A habilidade de ler as verdadeiras intenções de terceiros, identificando hesitações ou sinais de nervosismo.}
\pericia{Eloquência}{A arte da diplomacia e da persuasão, utilizada para convencer outros através de argumentos lógicos e carisma.}
\pericia{Enganação}{A capacidade de tecer mentiras convincentes, omitir fatos e manter disfarces sem levantar suspeitas.}
\pericia{Engenharia*}{O conhecimento técnico sobre a construção de estruturas, mecanismos complexos e fortificações.}
\pericia{Estratégia}{A capacidade de analisar o cenário em volta, identificar pontos cegos e planejar táticas de combate ou cerco.}
\pericia{Furtividade}{A perícia de mover-se silenciosamente e utilizar o ambiente para ocultar sua própria presença.}
\pericia{Intimidação}{O uso da presença física ou ameaças diretas para dobrar a vontade de outros através do medo.}
\pericia{Investigação}{A busca minuciosa por evidências e o gracioso loot, permitindo deduzir eventos passados através de pequenos detalhes ou encontrar itens específicos.}
\pericia{Luta}{A proficiência no combate corpo a corpo, englobando o uso de armas brancas e técnicas de combate desarmado.}
\pericia{Manutenção*}{A habilidade de reparar, afiar lâminas e garantir que o inventário permaneça funcional.}
\pericia{Medicina*}{O conhecimento de anatomia e farmacologia para tratar ferimentos, doenças e estabilizar aliados feridos gravemente.}
\pericia{Navegação*}{A capacidade de pilotar embarcações e dirigíveis.}
\pericia{Natureza}{O conhecimento sobre tudo que é natural e científico.}
\pericia{Percepção}{A agudeza dos sentidos para notar detalhes ambientais, sons distantes e presenças ocultas.}
\pericia{Performance}{A capacidade artística de impressionar ou distrair um público através de música, dança ou atuação.}
\pericia{Pontaria}{A precisão técnica com armas de longo alcance, como arcos, bestas e armas de fogo rúnicas.}
\pericia{Provocação}{A habilidade de insultar ou distrair oponentes, forçando-os a perder o foco tático ou atacar alvos específicos.}
\pericia{Reflexos}{A rapidez de resposta a estímulos súbitos, essencial para reagir a armadilhas e projéteis inesperados.}
\pericia{Sedução}{O uso do magnetismo pessoal e charme para influenciar as emoções e desejos de outros indivíduos.}
\pericia{Sobrevivência}{A aptidão para encontrar água, comida, rastrear presas e sobreviver em biomas hostis.}
\pericia{Temperança}{A resiliência mental e autocontrole, permitindo resistir ao pânico, tortura psicológica e corrupção mental.}
\pericia{Xilografia*}{A arte sagrada do \textit{RuneCrafting}, permitindo a gravação de runas e o trabalho em materiais condutores.}

\section*{Perícias e Progressão}

A maestria sobre as perícias em \textit{Overgrown} não é estática. Ele exige que o aventureiro escolha entre a versatilidade de aprender novos conceitos ou a profundidade de especializar-se naquilo que já conhece.

\begin{rpgboxtips}{Progressão de Perícia}
Diferente dos Atributos, as Perícias evoluem conforme o personagem interage com o Pináculo. Algumas habilidades na \textbf{Árvore de Talentos} podem conceder bônus diretos nestes valores ou permitir que você utilize uma perícia em situações não convencionais.
\end{rpgboxtips}

\subsection*{A Escolha de Aprendizado}

Todo personagem inicia sua jornada no \textbf{Nível 0} com \textbf{5 Perícias Iniciais} e 1 perícia de \link{data:origem}{Origem} (Maestria I). A cada \textbf{5 Níveis de Divindade} alcançados, o jogador deve escolher entre um dos dois caminhos:

\begin{enumerate}
    \item \textbf{Caminho da Versatilidade:} Aprender \textbf{duas novas perícias} em Maestria I (+5).
    \item \textbf{Caminho da Especialização:} Melhorar \textbf{uma perícia já existente} para o próximo nível de Maestria.
\end{enumerate}

\subsection*{Níveis de Maestria e Limites}

A evolução de uma perícia amplia o bônus de teste, mas o conhecimento avançado é trancado pela \link{att:divinity}{Divindade}. Você só pode realizar a melhoria se possuir o nível necessário:

\begin{center}
    \small 
    \renewcommand{\arraystretch}{1.4}
    \setlength{\tabcolsep}{4pt}
    \begin{tabularx}{\columnwidth}{@{} X r r @{}}
        \toprule
        \rowcolor{gray!15} 
        \textbf{\ringm{Maestria}} & \textbf{\ringm{Bônus}} & \textbf{\ringm{Divindade}} \\ 
        \midrule
        I \textit{(Iniciado)}     & +5  & 0+ \\
        II \textit{(Proficiente)} & +10 & 0 -- 19 \\
        III \textit{(Especialista)}& +15 & 20 -- 39 \\
        IV \textit{(Mestre)}      & +20 & 40 -- 50 \\
        \bottomrule
    \end{tabularx}
\end{center}

\begin{rpgboxtips}{Exemplo de Evolução}
No Nível 5, um jogador segue o \textbf{Caminho da Especialização} e sobe \textit{Luta} para +10. No Nível 10, ele escolhe \textbf{Versatilidade} e aprende \textit{Arcanismo} e \textit{Medicina} (+5). Ele só poderá levar \textit{Luta} para +15 ao atingir o Nível 20.
\end{rpgboxtips}